% !TeX root = ../nc-quadric-classification.tex
% 2026-11-21 / v0.1
\section{二次包絡の障害と二次超曲面表示}
\label{sec:quadric-envelope-obstruction}

以下、$k$ は代数閉体とし、代数は $k$ 上の局所有限な連結次数付き代数とする。
次数一で生成される代数 $B$ に対し、$V=B_1$ と置き、通常の二次包絡を
\[
 \operatorname{QEnv}(B)
 :=T(V)/\bigl(\ker(V\otimes V\longrightarrow B_2)\bigr)
\]
と定める。括弧は両側イデアルの生成を表す。
これは $B$ の二次関係をすべて課し、三次以上の独立な関係を忘れる構成である。
したがって $B$ 自身が二次代数ならば
$\operatorname{QEnv}(B)=B$ である。

型 $(1,4,6,4,1)$ は定義~\ref{def:as14641}の意味で用いる。

\begin{proposition}[次数二の正則元による商の数値]
\label{prop:quadric-hilbert}
$A$ を上記の型の代数、$g\in A_2$ を非零の両側正則正規元とし、
$B=A/(g)$ と置く。このとき $B$ は二次代数であり、
\[
 H_B(t)=\frac{1-t^2}{(1-t)^4}
       =\frac{1+t}{(1-t)^3},\qquad
 \dim_k B_n=(n+1)^2\quad(n\geq0).
\]
特に $\dim B_1=4$、$\dim B_2=9$ であり、
$\dim\ker(B_1\otimes B_1\to B_2)=7$ である。
\end{proposition}

\begin{proof}
正規性から $(g)=gA=Ag$ であり、左から $g$ を掛ける写像によって
右加群の完全列
\[
 0\longrightarrow A(-2)\xrightarrow{\,g\cdot\,}A
 \longrightarrow B\longrightarrow0
\]
を得る。したがって Hilbert 級数は主張の通りである。
係数は $\binom{n+3}{3}-\binom{n+1}{3}=(n+1)^2$ である
（第二項は $n<2$ なら零とする）。
$A$ が二次代数で $g$ が次数二なので、$B$ も二次代数である。
最後の次元は $4^2-9=7$ から従う。
\end{proof}

\begin{corollary}[通常の二次包絡を用いる定義の障害]
\label{cor:quadric-envelope-fails}
命題~\ref{prop:quadric-hilbert} の商 $B$ について
$\operatorname{QEnv}(B)=B$ であり、これは型 $(1,4,6,4,1)$ の
AS 正則代数ではない。
\end{corollary}

\begin{proof}
$B$ は二次代数なので二次包絡と一致する。
その次数二成分の次元は $9$ であるが、指定された型の AS 正則代数では
$10$ である。
\end{proof}

とくに $Q=\mathbb P^1\times\mathbb P^1$、
$\mathcal O=\mathcal O_Q$、$s=-\otimes\mathcal O_Q(1,1)$ に対しては、
$s^2\simeq S^{-1}[2]$ であり、$(\mathcal O,s)$ は Artin--Zhang ample である。
しかしその切断代数は
\[
 B=\bigoplus_{n\geq0}H^0(Q,\mathcal O_Q(n,n))
   \simeq k[x,y,u,v]/(xy-uv)
\]
なので、通常の二次包絡の AS 正則性を要求するとこの基本例を除外してしまう。
従って quadric における ambient 条件は、二次包絡の正則性から、
\emph{次数二の正則正規元を核に持つ AS 正則な二次被覆の存在}へ置き換える。
$s$ も被覆も、$\mathcal O$-marked object の指定データには含めず、
存在条件の証人として扱う。

\section{固定された切断代数の ambient 被覆}
\label{sec:ambient-incidence}

ここでは $B$ を二次代数とし、$V=B_1$、
$K=\ker(V\otimes V\to B_2)$ と置く。
さらに $\dim V=4$ および
$H_B(t)=(1-t^2)/(1-t)^4$ を仮定するので、$\dim K=7$ である。

\begin{lemma}[正規性と Hilbert 級数による正則性の判定]
\label{lem:normal-hilbert-regular}
$A$ は局所有限な連結次数付き代数、$0\ne g\in A_d$ は正規元とし、
$B=A/(g)$ とする。もし
$H_B(t)=(1-t^d)H_A(t)$ ならば、$g$ は左右の零因子でない。
\end{lemma}

\begin{proof}
右加群の完全列
\[
 0\longrightarrow\ker(g\cdot:A(-d)\to A)
 \longrightarrow A(-d)\xrightarrow{\,g\cdot\,}A
 \longrightarrow B\longrightarrow0
\]
を考える。正規性により像は商の核 $(g)$ に等しい。
Hilbert 級数を取ると、仮定は左端の核の Hilbert 級数が零であることを意味する。
局所有限性により核は零である。右から掛ける写像についても、
$Ag=(g)$ を用いて同じ議論を適用できる。
\end{proof}

\begin{proposition}[被覆の超平面表示]
\label{prop:ambient-hyperplane}
固定された $B$ に対する、次数一の同型 $A_1\to B_1=V$ を持つ
型 $(1,4,6,4,1)$ の AS 正則二次被覆
$A\twoheadrightarrow B$ で、その核が次数二の正則正規元で生成されるものは、
\[
 R\in\operatorname{Gr}(6,K)\simeq\mathbb P(K^*)
\]
のうち、次の二条件を満たすものによって分類される。
\begin{enumerate}
 \item $A_R:=T(V)/(R)$ は型 $(1,4,6,4,1)$ の AS 正則代数である。
 \item $(K/R)\subset(A_R)_2$ を張る任意の非零元 $g_R$ は正規元である。
\end{enumerate}
ここで被覆の同型は $B$ への写像と可換な次数付き代数同型とし、
核の生成元自体のスカラー倍は区別しない。
第二条件の下では $g_R$ の両側正則性は自動的に従う。
第一条件の下で、第二条件は二つの線形写像
\[
 V\longrightarrow(A_R)_3,\quad v\longmapsto g_Rv,
 \qquad
 V\longrightarrow(A_R)_3,\quad v\longmapsto vg_R
\]
がともに階数 $4$ を持つことと同値である。
\end{proposition}

\begin{proof}
与えられた被覆から
$R=\ker(V\otimes V\to A_2)$ を取ると、$\dim A_2=10$ より
$\dim R=6$ であり、合成写像から $R\subset K$ である。
二次性により $A=T(V)/(R)$ である。
核の次数二成分は $K/R$ であり、仮定された正則正規元の張る直線に等しい。

逆に $R$ が二条件を満たすとする。
$B=T(V)/(K)$ なので、$A_R/(g_R)=B$ である。
Hilbert 級数の仮定と第一条件から
$H_B(t)=(1-t^2)H_{A_R}(t)$ であり、
補題~\ref{lem:normal-hilbert-regular} によって $g_R$ は両側正則である。
この二構成は互いに逆である。実際、$B$ への写像と可換な被覆の同型は
次数一の同一視 $A_1=V$ を保存しなければならず、$A$ は次数一で生成されるので
同じ $R$ を与える。

階数による判定について、$J=\ker(A_R\to B)$ と置くと
$\dim J_3=\dim(A_R)_3-\dim B_3=20-16=4$ である。
二つの写像の像はともに $J_3$ に含まれるので、ともに階数 $4$ ならば
$g_RV=J_3=Vg_R$ である。次数一生成性を用いて生成元を一つずつ動かすと
$g_RA_R=A_Rg_R$ が従う。
逆に正規性が成り立てば、既に示した両側正則性から両写像は単射であり、
その階数は $4$ である。
\end{proof}

\begin{remark}[有限次数の判定と未解決部分]
$A_R$ は次数一で生成されるため、
$g_RV=Vg_R$ が $(A_R)_3$ の部分空間として成立することと
$g_R$ の正規性は同値である。
一方の包含から生成元を一つずつ動かして $g_RA_R\subset A_Rg_R$ を得て、
逆の包含も同様に示せるからである。
これは正規性の有限次元線形代数による判定であるが、
有限次数での Hilbert 関数の一致だけで AS 正則性まで証明できるとは主張しない。
また、この許容集合が開集合である、滑らかである、あるいは一点であるという
主張はここには含まれない。
\end{remark}

\section{同じ可換 quadric を持つ非同一の ambient 被覆}
\label{sec:ore-ambient-family}

\begin{proposition}[正規被覆を忘れる写像の非自明なファイバー]
\label{prop:ore-ambient-family}
$a\in k^{\times}$ に対し
\[
 A_a=k\langle x,y,u,v\rangle/
 \bigl([u,v],[u,x],[u,y],[v,x],[v,y],
       yx-a\,xy-(1-a)uv\bigr),\qquad g=xy-uv
\]
と定める。このとき次が成り立つ。
\begin{enumerate}
 \item $A_a$ は noetherian な二次 AS 正則 domain であり、型は
       $(1,4,6,4,1)$ である。
 \item $g$ は次数二の正則正規元であり、
       $gx=a xg$、$gy=a^{-1}yg$、$gu=ug$、$gv=vg$ を満たす。
 \item 生成元を同一視する標準同型
       \[
        A_a/(g)\simeq k[x,y,u,v]/(xy-uv)=:B
       \]
       がある。
 \item $B_1$ の同一視を固定した被覆としては、$a\ne b$ ならば
       $A_a\twoheadrightarrow B$ と $A_b\twoheadrightarrow B$ は同型でない。
       $A_1$ は可換であり、$a\ne1$ ならば $A_a$ は非可換で $g$ は非中心である。
\end{enumerate}
\end{proposition}

\begin{proof}
$R_0=k[u,v,x]$ とし、
\[
 \sigma(u)=u,\quad\sigma(v)=v,\quad\sigma(x)=ax,
 \qquad
 \delta(u)=\delta(v)=0,\quad\delta(x)=(1-a)uv
\]
と置く。$a\ne0$ なので $\sigma$ は次数を保つ自己同型である。
$\delta$ は次数一の $\sigma$-derivation に延びる。
たとえば可換関係 $xu-ux$ に対して
\[
 \delta(xu-ux)=(1-a)uvu-u(1-a)uv=0
\]
であり、$xv-vx$、$uv-vu$ についても同様である。
従って表示から
$A_a=R_0[y;\sigma,\delta]$ である。

この Ore 拡大は、多項式環の自明な Zhang twist を基礎環とする
次数一の Ore 拡大であるため、
\cite[Proposition~2.2]{LV2026}（$n=3$、twist は恒等写像）の仮定を満たす。
同命題により $A_a$ は noetherian AS 正則 domain、global dimension $4$、
Koszul である。この適用に必要な条件は直前で確認した。
Ore 拡大の基底
\[
 \{u^i v^j x^r y^s\mid i,j,r,s\geq0\}
\]
から $H_{A_a}(t)=(1-t)^{-4}$ を得る。
Koszul 性により極小自由分解の第 $i$ 項は次数 $i$ で生成されるので、
Hilbert 級数の逆数
\[
 (1-t)^4=1-4t+6t^2-4t^3+t^4
\]
から分解の型が $(1,4,6,4,1)$ であると分かる。

基底より $g=xy-uv\ne0$ である。$A_a$ は domain なので $g$ は両側正則である。
$u,v$ の中心性と表示関係から
\[
 \begin{split}
 gx&=xyx-uvx
     =x\bigl(a\,xy+(1-a)uv\bigr)-xuv
     =a x(xy-uv)=a xg,\\
 yg&=yxy-yuv
     =\bigl(a\,xy+(1-a)uv\bigr)y-uvy
     =a(xy-uv)y=a gy
 \end{split}
\]
を得る。これは主張された正規性の関係を与える。

商で $g=0$ を課すと $xy=uv$ であり、表示関係は
$yx=a uv+(1-a)uv=uv=xy$ となる。
残る関係はすべて可換関係なので、商はちょうど上記の $B$ である。

最後に、$W\subset V\otimes V$ を表示中の最初の五つの可換関係の張る部分空間、
$r_a=y\otimes x-a x\otimes y-(1-a)u\otimes v$ と置く。
$R_a=W+kr_a$ が被覆の二次関係空間である。
もし $R_a=R_b$ なら、$y\otimes x$ の係数を比較して
$r_a-r_b=(b-a)(x\otimes y-u\otimes v)\in W$ となる。
右辺の括弧内の元は $W$ に含まれないので $a=b$ である。
命題~\ref{prop:ambient-hyperplane} の同型に関する議論により、
これは固定された $B$ 上の被覆の非同型性を意味する。
$a\ne1$ なら $[y,x]=(a-1)g\ne0$ であり、
$gx-xg=(a-1)xg\ne0$ なので、非可換性と非中心性も従う。
\end{proof}

\begin{remark}[この命題が分類について述べる範囲]
上の族は、固定された切断代数 $B$ から ambient 正規被覆が一意に復元される
という主張への反例である。同じ
$(\operatorname{coh}(\mathbb P^1\times\mathbb P^1),\mathcal O,s)$ を与えるので、
異なる $a$ を異なる $\mathcal O$-marked quadric と数えてはならない。
ただし、基底の代数 $B$ の自己同型まで許す場合の $a$ の同値関係を
この命題は分類していない。次数二の中心元だけを許す被覆という別の問題も、
上記の正規被覆の問題と区別する必要がある。
この具体例の文献上の新規性は未確認であり、新分類定理として扱わない。
\end{remark}

